\scanpage{36}
% archive sheet number: 35
\begingroup\makeatletter\tagsleft@true\makeatother
if \deleted{\unsure{Card|the}} cardinality of the set $(18)$ is $4$, i.e. if the four subsets of $X$
\[
\Sigma_1^+,\quad\Sigma_1^-,\quad\Sigma_2^+,\quad\Sigma_2^-
\]
are distinct. In this connection let us note that we already know by \deleted{hypothesis $(\unsure{\Pi|\pi}_1,\Pi_2$ \unsure{being}} \inserted{definition} of binary déc.
\[
\Sigma_1^+\ne\Sigma_1^-,\qquad\Sigma_2^+\ne\Sigma_2^-
\tag*{($*$)}
\]
and that we have
\[
\Sigma_1^+\ne\Sigma_2^+
\tag*{($**$)}
\]
since $\Sigma_1^+$ and $\unsure{\Sigma|\Sigma_2}^{\unsure{+|\top}}$ are disjoint \uline{non} \uline{empty} subsets. Let us therefore note that we read from the diagram $(14)$
\[
\left\{\begin{aligned}
\text{a)}\;(\Sigma_1^+\unsure{\ne|=}\Sigma_2^-)&\Longleftrightarrow E_1\unsure{\ne|=}\Sigma_{12}
   &(\Longrightarrow E_2=\varnothing)\\
\text{b)}\;(\Sigma_2^+=\Sigma_1^-)&\unsure{\Longleftrightarrow|\Longrightarrow} E_2=\Sigma_{12}
   &(\Longrightarrow E_1=\varnothing)
\end{aligned}\right.
\tag*{($***$)}
\]
and that the two relations \inserted{\unsure{-|a}), b)} considered in $(***)$ cannot be satisfied simultaneously, since $\deleted{\Sigma_{12}}$ we would have $\Sigma_{12}=E_1=E_2=\varnothing$ since $E_1\cap E_2=\varnothing$, which contradicts $(16)$ (c).

Finally
\[
(\Sigma_1^-=\Sigma_2^-)\Longleftrightarrow
(\Sigma_{1\unsure{\deleted{2}|2}}^-=\unsure{\deleted{X|E}}\Sigma_2^-=X)
\quad\text{i.e. }\Sigma_{12}=X
\tag*{($****$)}
\]
This case is mutually exclusive with each of the cases $(***)$, \deleted{\unsure{in the latter}} \inserted{\unsure{For|for} if we had at the same time} $\deleted{\unsure{\Sigma_1^-}=\Sigma_2^-=X}\ \unsure{\deleted{E}}\Sigma_{12}=X$ and $E_1=\Sigma_{12}$, we would have $E_1=X$ hence $\Sigma_1^+=\Sigma_1^-=X$, \deleted{\unsure{of|and}}\unsure{which} contradicts $(3)$.

\uline{Proposition} For us to have $\operatorname{Card}\{\Sigma_1^+,\Sigma_1^-,\Sigma_2^+,\Sigma_{\unsure{2|1}}^-\}=4$, it is necessary and sufficient that we not be one of the following three cases, which exclude one another
\[
\left\{\begin{aligned}
\text{\unsure{a})}\quad &E_1=\Sigma_{12}&&(\Longrightarrow E_2=\varnothing)\\
\text{a')}\quad &E_2=\Sigma_{12}&&(\Longrightarrow E_1=\varnothing)\\
\text{b)}\quad &\Sigma_{12}=X\quad\text{i.e. }\Sigma_1^-=\Sigma_2^-=X
   &&(\text{whence }E_1=\Sigma_1^+,E_2=\Sigma_2^+)
\end{aligned}\right.
\tag*{(\deleted{\unsure{19|20}}20)}
\]
in which case the cardinality \unsure{of the|of the|of} $\unsure{\}\,4|\{\ \}|4}$ is equal \unsure{:|to} \unsure{three}.

\endgroup
\clearpage
